\section{Theorem and Provenance Ledger}

This section records the exact status of the structures used in
Paper 14.

The purpose is to keep imported theorem-level facts, new results,
internal conjectures, computational frontiers, and physical bridge
hypotheses distinct.

\subsection{Imported theorem-level structure}

The following objects are imported from the established finite
AT4val[60,6] theorem chain.

\begin{center}
\begin{tabular}{p{0.34\linewidth}p{0.56\linewidth}}
\toprule
Object or statement & Status \\
\midrule

\(
G_{15}\cong L(\mathrm{Petersen})
\)
&
proved finite graph identification
\\

\(
M\in\{0,1\}^{15\times30}
\)
&
proved transport-induced sector-edge incidence matrix
\\

row weight \(14\), column weight \(7\)
&
proved incidence census
\\

\(
Q=MM^{\mathsf T}
\)
&
proved sector-overlap Gram identity
\\

\(
Q=A^3+2A^2+2I
\)
&
proved adjacency-algebra identity
\\

\(
\operatorname{spec}(Q)
=
\{98^{(1)},18^{(5)},3^{(4)},2^{(5)}\}
\)
&
proved spectral decomposition
\\

signed
\(
G_{30}\to G_{15}
\)
lift
&
proved finite covering structure
\\

nontrivial binary switching class
&
proved lifted holonomy not determined by \(Q\)
\\

projected return before complete return
&
proved in multiple declared finite constructions
\\

registered history
&
proved retained finite distinction relative to declared projections
\\

\(
\tau_{\mathrm T}
=
\operatorname{wind}_{r_1}
\)
&
proved intrinsic additive oriented history coordinate
\\

\(
R_n(h)
=
\|h-\overline h\,\mathbf 1_n\|_2
\)
&
proved centered Euclidean invariant on a declared finite register
\\

\bottomrule
\end{tabular}
\end{center}

\subsection{New theorem-level results of Paper 14}

Paper 14 introduces the canonical normalization

\[
\widehat Q
=
\frac{1}{98}Q.
\]

The following results are proved in this paper.

\begin{center}
\begin{tabular}{p{0.42\linewidth}p{0.48\linewidth}}
\toprule
Statement & Status \\
\midrule

\(
\widehat Q\mathbf 1
=
\mathbf 1
\)
&
proved
\\

common-mode eigenvalue
\(
1
\)
&
proved
\\

centered eigenvalues
\(
9/49,\ 3/98,\ 1/49
\)
&
proved
\\

centered spectral radius
\[
\rho_{\mathrm{cent}}
=
\frac{9}{49}
\]
&
proved
\\

\[
R(\widehat Q^{\,k}x)
\le
\left(\frac{9}{49}\right)^kR(x)
\]
&
proved
\\

\[
\widehat Q^{\,k}
\longrightarrow
P_0
\]
&
proved in Euclidean operator norm
\\

\[
\left\|
\widehat Q^{\,k}-P_0
\right\|_{2\to2}
=
\left(\frac{9}{49}\right)^k
\]
&
proved and sharp
\\

\(
\det\widehat Q\neq0
\)
&
proved
\\

finite-stage invertibility of
\(
\widehat Q^{\,k}
\)
&
proved for every finite \(k\)
\\

rank-one asymptotic limit
&
proved because
\(
\operatorname{rank}(P_0)=1
\)
\\

finite-resolution contrast loss
&
proved relative to any fixed positive threshold
\\

\bottomrule
\end{tabular}
\end{center}

\subsection{Exact theorem keeper}

The new theorem package may be summarized as:

\begin{quote}
The normalized Thalean quadratic is an invertible finite linear
contrast concentrator whose iterates converge sharply to the
common-mode projection.
\end{quote}

Equivalently:

\begin{quote}
Relative registered contrast may approach zero while exact
finite-stage linear distinction remains preserved.
\end{quote}

\subsection{Defined but not promoted to theorem}

The following terms are defined in Paper 14 but are not themselves
proved native dynamical laws:

\begin{itemize}

\item projection-relative saturation;

\item contrast saturation;

\item effective registration degeneracy;

\item clean release;

\item overcompression;

\item inter-face sigma;

\item admitted summation domain.

\end{itemize}

These definitions organize the conjectural mechanics around the exact
linear theorem.

\subsection{Thalean Relational Saturation Conjecture}

The central conjectural extension is:

\begin{quote}
Given a globally compatible finite linear admissible transport
structure, a declared projection or registration, a finite additive
aggregation law, and finite independent distinction capacity,
increasing relational loading can strengthen the coherent aggregate
of bounded admissible contributions while reducing their independent
registered contrast.

Beyond a critical regime, additional loading is expressed through
increasing dependency, common-mode concentration, and
projection-relative degeneracy rather than unbounded growth of
independent local distinction.

Complete distinctions may remain preserved by the finite transport
while the declared registration no longer resolves them independently.
\end{quote}

This statement is not yet a theorem of the native higher mechanics.

\subsection{Missing native crosswalk}

The principal missing mathematical arrow is

\[
\text{native admitted transport/loading}
\longrightarrow
\text{certified contrast-concentrating register}.
\]

A strong theorem would require a diagram of the form

\[
\begin{array}{ccc}
X_k & \xrightarrow{T} & X_{k+1} \\
\downarrow \rho & & \downarrow \rho \\
H_k & \xrightarrow{L} & H_{k+1},
\end{array}
\]

with

\[
\rho\circ T
=
L\circ\rho,
\]

where:

\begin{enumerate}[label=(\roman*)]

\item \(T\) is a native admitted transformation;

\item \(\rho\) is a certified registered readout;

\item \(L\) is contrast-concentrating;

\item retained complete distinction remains nontrivial upstairs.

\end{enumerate}

Paper 14 proves the required behavior for the candidate operator

\[
L=\widehat Q.
\]

The native \(T\) and \(\rho\) crosswalk remain open.

\subsection{Current inter-face sigma frontier}

The completed Project 41 Mac search used the canonically completed,
neighbor-frame-\(E\)-closed local seam vocabulary containing

\[
2816
\]

local options.

The completed search reached:

\begin{itemize}

\item \(376\) total solver iterations;

\item \(175069\) unique four-step closure nogoods;

\item repeated best candidates with
\[
1888
\]
failed flags out of
\[
1920;
\]

\item equivalently,
\[
32
\]
correctly closing flags;

\item no complete global \(r_3\);

\item no proof that the completed search space is infeasible.

\end{itemize}

The exact status is therefore

\[
\text{deeper bounded partial closure}
\]

but not

\[
\text{global transport theorem}.
\]

\subsection{Sigma boundary}

The unresolved inter-face sigma is not identified with:

\begin{itemize}

\item the signed \(G_{30}\to G_{15}\) cocycle;

\item the twelve-section binary closure class;

\item either integer parent;

\item the primitive registered-history cycle \(r_1\);

\item Thalean time;

\item or the normalized quadratic \(\widehat Q\).

\end{itemize}

An explicit structural crosswalk is required before any such
identification.

\subsection{Compatibility versus saturation}

The following distinction is locked:

\[
\text{local inadmissibility}
\neq
\text{global incompatibility}
\neq
\text{relational saturation}.
\]

In particular,

\[
\text{failure to find global sigma}
\]

is not evidence for

\[
\text{saturation}.
\]

The saturation conjecture applies only after a globally compatible
transport structure has been admitted.

\subsection{Finite summation ledger}

Paper 14 uses three established forms of finite accumulation.

\begin{center}
\begin{tabular}{p{0.27\linewidth}p{0.32\linewidth}p{0.31\linewidth}}
\toprule
Structure & Law & Domain \\
\midrule

incidence sum
&
\(
x\mapsto MM^{\mathsf T}x
\)
&
fifteen-sector real register
\\

binary regional sum
&
\(
d\delta=\omega
\)
&
twelve-section binary closure complex
\\

integer regional sum
&
\(
d\eta=\Omega
\)
&
integer refinement of closure complex
\\

history winding
&
\(
\tau_{\mathrm T}(\gamma_2\circ\gamma_1)
=
\tau_{\mathrm T}(\gamma_2)
+
\tau_{\mathrm T}(\gamma_1)
\)
&
registered-history cycle complex
\\

\bottomrule
\end{tabular}
\end{center}

These are related examples of lawful finite accumulation.

They are not one universal sum.

\subsection{Clean-release status}

The clean-release sequence

\[
\text{organization}
\to
\text{clean release}
\to
\text{overcompression}
\to
\text{contrast saturation}
\]

is an internal hypothesis.

The quadratic theorem proves only the contrast-concentrating direction.

No theorem in Paper 14 proves:

\begin{itemize}

\item a unique clean-release optimum;

\item a native loading parameter;

\item a universal release statistic;

\item or a physical emission mechanism.

\end{itemize}

\subsection{Empirical ledger}

The existing WMAP \(C_3/\Xi\) experiment is retained with its frozen
decision.

The exact empirical status is:

\[
\text{all four primary cells in predicted upper tail},
\]

\[
\text{three of four below }p=0.05,
\]

\[
p_{\mathrm{W,Gaussian}}
=
0.0539461,
\]

and therefore

\[
\boxed{\text{NO-GO REPLICATION}}.
\]

The experiment is not reinterpreted as a relational-saturation test.

Any saturation experiment requires a new frozen protocol.

\subsection{Physical bridge status}

The following remain physical bridge hypotheses rather than theorem
consequences:

\begin{itemize}

\item finite physical independent-distinction density;

\item minimum four-support of order \(\ell_P^4\);

\item physical gravity;

\item black-hole relational saturation;

\item horizon-like outward registration loss;

\item singularity avoidance;

\item black-hole entropy;

\item holography;

\item Hawking radiation;

\item a physical CMB saturation mechanism;

\item physical proper time;

\item Standard Model correspondence.

\end{itemize}

\subsection{Explicitly rejected identifications}

Paper 14 explicitly rejects the unsupported implications

\[
Q
=
\text{native temporal evolution},
\]

\[
\widehat Q
=
\text{physical dynamics},
\]

\[
\text{common mode}
=
\text{physical mass},
\]

\[
\text{centered magnitude}
=
\text{physical information},
\]

\[
\text{contrast concentration}
=
\text{gravity},
\]

\[
\text{projection-relative saturation}
=
\text{event horizon},
\]

\[
\text{finite saturation}
=
\text{singularity resolution},
\]

and

\[
\Xi
=
\text{saturation observable}.
\]

\subsection{Current theorem frontier}

The final status of Paper 14 is:

\[
\text{quadratic incidence law}
\quad
\text{proved},
\]

\[
\text{contrast concentration}
\quad
\text{proved},
\]

\[
\text{finite-stage non-erasure}
\quad
\text{proved},
\]

\[
\text{projection-relative saturation}
\quad
\text{defined},
\]

\[
\text{native loading crosswalk}
\quad
\text{open},
\]

\[
\text{global inter-face sigma}
\quad
\text{open},
\]

\[
\text{native relational-saturation theorem}
\quad
\text{open},
\]

\[
\text{physical saturation bridge}
\quad
\text{open}.
\]

\subsection{Final keeper}

The theorem frontier can be stated in one sentence:

\begin{quote}
Paper 14 proves how a finite linear register can lose relative
individuation without finite-stage state merger; the remaining problem
is to prove that an admitted native Thalean transport actually enters
that regime while retained complete history remains nontrivial.
\end{quote}
